\begin{abstract}
    Le password restano il meccanismo di autenticazione più diffuso, ma la loro sicurezza dipende in larga misura dalle abitudini degli utenti, che tendono a sceglierle a partire da informazioni personali facili da ricordare.
    L'Open Source Intelligence (OSINT), ovvero la raccolta e l'analisi di informazioni pubblicamente disponibili su social network, siti web e altre fonti aperte, consente di ricostruire con poco sforzo nomi, date, interessi e relazioni di un individuo.
    Si esamina come questi dati possano essere integrati negli attacchi a dizionario, generando liste di parole mirate che aumentano sensibilmente la probabilità di indovinare una password rispetto ai dizionari generici.
    Si analizzano inoltre i fattori che rendono le password prevedibili e si discutono le contromisure, sia tecniche sia comportamentali, per ridurre l'efficacia di questi attacchi.
    Il lavoro evidenzia come la sicurezza delle credenziali sia strettamente legata alla quantità di informazioni che gli utenti espongono online.
\end{abstract}
